\subsection{The Selmer argument and its missing constant}

Throughout this subsection, a section means a continuous splitting modulo
conjugation by the geometric fundamental group. We distinguish the classical
curve from the plus curve. The main conclusion of LM6 E.6--E.12 survives the
following audit, but one of its local assertions is false: a ramification index
does not remove the Kummer class of the leading coefficient. We give the
correction before invoking residual modularity.

\begin{statement}{A.1}{Selmer sections and passage to a neighbourhood}{Known}
Let $U/k$ be a smooth geometrically connected hyperbolic curve over a number
field, with smooth compactification $X$ and reduced boundary $D$.
A Selmer section is one whose restriction at every completion $k_v$ is the
section of a point of $U(k_v)$ or belongs to the packet of a cusp in $D(k_v)$.
There are no exclusions at $2$ or at a chosen coefficient prime.
A finite \'etale neighbourhood of a Selmer section again carries a Selmer
section. Its compactification gives a modified adelic point surviving all
finite descent.
\end{statement}
\begin{proof}
The definition is that of \citet[\S3.3]{AMSteklov2026}. At a real place a
point section specifies the connected component, and not a distinguished
point of that component. If $H\subset\pi_1(U)$ is the open subgroup defining
a neighbourhood, containment of a local point section in a conjugate of $H$
is precisely a fixed coset in the finite geometric fibre. It therefore gives
a rational point of that fibre. At a cusp we use the corresponding finite
fibre of the normalization of $X$. Intersecting the cusp decomposition group
with the conjugate of $H$ gives the decomposition group upstairs; since the
given section surjects onto $G_{k_v}$, the cusp upstairs has residue field
$k_v$. Thus the lifted section has the required local description.

After compactification both kinds of local sections become sections of
local points. The hypotheses of \citet[Theorem 2.1]{AMHarariStix2012}
are now satisfied with the geometric quotient equal to the full geometric
fundamental group and with no omitted places. Their implication
(iii)$\Rightarrow$(i) says that these local points survive every finite
torsor. Smoothness, properness and geometric connectedness hold for the
compactification of the connected neighbourhood: the section forces its
constant field to be $k$. This proves the asserted finite-descent statement.
\end{proof}

\begin{statement}{A.2}{Finite support, including the real-place correction}{Proved}
Let $U/\mathbb Q$ and $X$ be as above, let $s$ be Selmer, and let
$Z\subset X$ be a finite reduced subscheme. If the underlying local points
of $s$ lie in $Z(\mathbb Q_p)$ for a density-one set of finite primes, then
$s$ is a point or cusp section above an element of $Z(\mathbb Q)$.
Consequently a nonconstant $\mathbb Q$-morphism $f:X\to W$ to a proper curve
with
\[
 W(\mathbb A_{\mathbb Q})_\bullet^{\mathrm{f\text{-}cov}}=W(\mathbb Q)
\]
makes every Selmer section of $U$ point-induced or cuspidal.
\end{statement}
\begin{proof}
We may choose a cofinal family of neighbourhoods $U'\to U$ whose smooth
compactifications $X'$ have genus at least one. Here is an explicit
justification. After any given neighbourhood, a compactification of positive
genus needs no refinement for this purpose. Otherwise it has genus zero
and at least $r=3$ punctures. Take the kernel of geometric mod-$3$
abelianization. This characteristic subgroup, together with the section,
defines a neighbourhood over $\mathbb Q$. Its degree is $d=3^{r-1}$ and
each boundary inertia group has ramification index three. Hence
\[
 2g(X')-2=d\left(-2+\frac{2r}{3}\right)\geq0.
\]
This construction after an arbitrary initially prescribed neighbourhood
proves cofinality.

Put $Z'=(X'\times_X Z)_{\mathrm{red}}$. By A.1 the lifted section gives
an adelic point of $X'$ surviving finite descent, and therefore finite
abelian descent. It belongs to $Z'(\mathbb Q_p)$ on the same density-one
set. The hypotheses of \citet[Theorem 8.2]{AMStoll2007} hold: $X'$ is a
smooth proper curve of genus at least one, and $Z'$ is finite. The theorem
gives a point of $Z'(\mathbb Q)$.

These finite nonempty sets of rational points have nonempty inverse limit.
Indeed, a common refinement of any finite collection of neighbourhoods
supplies a point mapping compatibly to all of them. The resulting pro-point
over $z\in Z(\mathbb Q)$ is fixed by $s(G_{\mathbb Q})$. Its stabilizer
is a point decomposition group if $z\in U$, and a cusp decomposition group
otherwise. One can see this directly on the inverse limits of geometric
fibres: they are the homogeneous spaces of the geometric fundamental group
modulo the trivial inertia group at an interior point, or modulo the cusp
inertia group at the boundary. Thus $s$ belongs to the asserted packet.

For the consequence, functoriality of finite descent sends the adelic
localization to one diagonal $w\in W(\mathbb Q)$. The fibre $f^{-1}(w)$ is
finite because $f$ is a nonconstant morphism of proper smooth curves.
Apply the first assertion. No separate hypothesis that all of $W(\mathbb Q)$
is finite is needed for this implication.

We use finite places only in forming the fibre support. This is exactly the
correction to Proposition 8.5 in \citet[\S3]{AMStollErrata2017}: a real
connected component mapping to the component of $w$ need not contain a real
point in the scheme-theoretic fibre above $w$. The density-one formulation
of Theorem 8.2 removes that problem. The erratum's separate warning about
components of abelian torsors does not license a shortcut here; our adelic
point survives \emph{all} finite torsors by A.1.
\end{proof}

\begin{statement}{A.3}{The affine map on cusp packets}{Proved}
Let $k$ have characteristic zero. Suppose a morphism of completed punctured
local curves has equation $q=a u^e h(u)$, where $a\in k^\times$,
$h\in1+u k[[u]]$ and $e>0$. Relative to the origins given by $q$ and $u$,
the induced map on section packets is
\[
 H^1(k,\widehat{\mathbb Z}(1))\longrightarrow
 H^1(k,\widehat{\mathbb Z}(1)),\qquad
 b\longmapsto\kappa(a)+eb.
\]
For the Legendre family on $U_\ell$, at every rational local cusp and every
prime $\ell\geq5$, the Tate parameter after removing its quadratic
splitting character satisfies
\[
 q(s)=-8\kappa(2)+2\ell\,s^*\kappa(u).
 \tag{A3.1}
\]
\end{statement}
\begin{proof}
The group $1+u k[[u]]$ is uniquely divisible: its $m$th roots with constant
term one exist uniquely by Hensel's lemma for every positive integer $m$,
and are compatible. Thus its Kummer class vanishes. Pulling the Kummer
class of $q$ back along a section gives the first formula.

For Fermat, $t=-X^\ell/Z^\ell$ and
\[
 j(t)=2^8\frac{(t^2-t+1)^3}{t^2(t-1)^2}.
\]
At $X=0$, $Y=0$ and $Z=0$ respectively use $u=X/Z$, $u=Y/Z$ and
$u=Z/X$. The three formulas for $t$ are
\[
 -u^\ell,\qquad 1+u^\ell,\qquad -u^{-\ell}.
\]
In each case $j^{-1}=2^{-8}u^{2\ell}h_0(u)$ with $h_0(0)=1$.
The formal inverse of the Tate expansion $j=q^{-1}+744+\cdots$ has
$q=j^{-1}+O(j^{-2})$. Hence $q=2^{-8}u^{2\ell}h(u)$ with $h(0)=1$,
and the preceding argument proves (A3.1). These coordinates work at every
cusp in the indicated divisor, including cusps rational only over $k$.
\end{proof}

\begin{statement}{A.4}{A false line in LM6 E.8 and its source}{Refuted}
The assertion in LM6 E.8 that the cusp Tate parameter is divisible by
$\ell$, and the general ramification-divisibility assertion in
\citet[Proposition labelled \texttt{prop:ramificationindexforcuspidal}]{AMSteklov2026},
are false. The correct statement is A.3.
\end{statement}
\begin{proof}
Take $k=\mathbb Q_2$, $\ell=5$, the rational cusp $X=0$ of $U_5$, and the
tangential section with $s^*\kappa(u)=0$. Equation (A3.1) gives
$q(s)=-8\kappa(2)$. Its image in $H^1(\mathbb Q_2,\mu_5)$ has valuation
$-8\equiv2\pmod5$, so it is nonzero. Thus the parameter is not divisible
by five and the residual Tate extension is nonsplit. A finite witness is
\[
 \tau=\begin{pmatrix}1&2\\0&1\end{pmatrix},\qquad
 \varphi=\begin{pmatrix}2&0\\0&1\end{pmatrix}
 \quad\text{in }\mathrm{GL}_2(\mathbb F_5).
\]
They satisfy $\tau^5=1$, $\tau\ne1$ and
$\varphi\tau\varphi^{-1}=\tau^2$, and generate a group of order twenty.
The exact program \texttt{a\_main\_exact.py} checks these assertions and
the three leading coefficients. The erroneous source replaces a local
map of ramification index $e$ by $q=u^e$ over $k$, thereby dropping the
coefficient $a$. Such a replacement need not exist over $k$; over an
algebraic closure it loses the arithmetic datum under discussion.
\end{proof}

\begin{statement}{A.5}{The corrected local conductor and weight calculation}{Proved}
Let $\ell\geq5$ and let a local section of $U_\ell/\mathbb Q_p$ be
point-induced or cuspidal. There is a quadratic character $\delta_p$ such
that its twisted mod-$\ell$ Legendre representation is unramified when
$p\notin\{2,\ell\}$, finite flat of Serre weight two when $p=\ell$, and
tamely ramified of conductor exponent one when $p=2$. The splitting
character is allowed to be ramified. A twist with the same inertial
restriction has the same stated properties.
\end{statement}
\begin{proof}
For a point, write $a^\ell+b^\ell+1=0$ and $t=-a^\ell$.
If $v_p(a)\ne0$, direct evaluation of $j$ gives
\[
 v_p(j)=8v_p(2)-2\ell|v_p(a)|.
 \tag{A5.1}
\]
If $v_p(a)=0$ and $t\equiv1\pmod p$, use $1-t=-b^\ell$ instead.
Otherwise $p$ is odd and the integral Legendre equation has unit
discriminant $16t^2(1-t)^2$, so it has good reduction. At $2$ this last
case is impossible, and (A5.1) is negative because $\ell\geq5$.
Thus at odd primes we have good or potentially multiplicative reduction,
and at $2$ we have potentially multiplicative reduction.

Remove the quadratic splitting character in the latter case. The result
is a Tate representation, an extension of $\mathbb F_\ell$ by
$\mathbb F_\ell(1)$ with Kummer parameter $q$ and
$v_p(q)=-v_p(j)$. For odd $p\ne\ell$ this valuation is divisible by
$\ell$. Units acquire $\ell$th roots over the maximal unramified extension,
so the residual extension is unramified. Good reduction gives the same
conclusion. At $p=\ell$ good reduction supplies the finite flat model
$E[\ell]$; in the Tate case the divisibility of $v_\ell(q)$ gives a unit
representative modulo $\ell$th powers. Its Kummer extension extends to a
finite flat extension of $\mathbb Z/\ell$ by $\mu_\ell$ over
$\mathbb Z_\ell$. The weight is two by \citet[Propositions 4--5,
\S\S2.8--2.9]{AMSerre1987}. These statements include both ordinary and
supersingular good reduction.

At $2$ we have $v_2(q)=2\ell m-8\not\equiv0\pmod\ell$. Wild inertia
is pro-$2$ and cannot map nontrivially to the additive group of
$\mathbb F_\ell$. The nontrivial tame unipotent inertia fixes exactly one
line; the Artin conductor exponent is consequently one.

For a cusp, A.3 gives the residual extension class
$-8\kappa_\ell(2)$. It is an unramified class at $p\ne2,\ell$, a unit
Kummer class at $p=\ell$, and a nonzero valuation class at $2$. The same
arguments apply. In particular nonsplitting at a cusp does not damage
finite flatness at $\ell$.

Finally, two splitting twists with equal restriction to inertia differ by
an unramified character. Unramified twisting preserves the conductor and
preserves finite flatness by finite \'etale descent. This explains explicitly
why ramified splitting characters are harmless once they are removed;
one must not assume that the original representation already has the
twisted conductor or flatness properties.
\end{proof}

\begin{statement}{A.6}{Residual reducibility and the two characters}{Known}
For a Selmer section of $U_\ell/\mathbb Q$, $\ell\geq5$, the mod-$\ell$
Legendre representation is reducible. A global quadratic twist has
semisimplification $1\oplus\overline\chi_\ell$. Quadratic twisting preserves
the trivial mod-two representation. In the irreducibility argument the
twisted Serre level is in fact $2$, stronger than the bound $16$ in LM6 E.8.
\end{statement}
\begin{proof}
A continuous representation into a finite group is ramified at only
finitely many primes. The splitting characters of A.5 can be chosen
unramified outside a finite set: at a good prime we choose the trivial
character; outside $2\ell$, a ramified scalar quadratic character could
not turn an unramified residual representation into another unramified
one when $\ell$ is odd. The cyclotomic reciprocity isomorphism
$G_{\mathbb Q}^{\mathrm{ab}}\simeq\prod_p\mathbb Z_p^\times$ patches
their inertia restrictions into a global quadratic character $\delta$.
The quotient of its localization by each chosen character is unramified.
By A.5 the resulting representation has conductor $2$ and weight two;
its determinant is still $\overline\chi_\ell$.

Suppose it were irreducible. It is absolutely irreducible: complex
conjugation has the distinct eigenvalues $1,-1$, so any invariant line
over $\overline{\mathbb F}_\ell$ would be one of two lines already
defined over $\mathbb F_\ell$. It is odd by its determinant. The precise
Serre modularity theorem therefore gives a weight-two eigenform of
level two. Its character is trivial since
$(\mathbb Z/2)^\times$ is trivial. But $X_0(2)$ has genus zero, so there
is no such cusp form. We use the general-conductor theorem of
\citet[Theorem 9.1]{AMKhareWintenbergerI2009}, with hypothesis (H)
supplied by \citet[Theorem (0.1) and Corollary (0.2)]{AMKisin2009};
Theorem 1.2(1) of that paper alone has an odd-conductor hypothesis and
does not suffice. The companion lifting results are in
\citet{AMKhareWintenbergerII2009}. For comparison with the weaker
printed bound, $X_0(N)$ has genus zero for each $N=1,2,4,8,16$.
The respective tuples
$(\text{index},\text{cusps},e_2,e_3)$ are
$(1,1,1,1)$, $(3,2,1,0)$, $(6,3,0,0)$, $(12,4,0,0)$, and
$(24,6,0,0)$; substitution in
$1+\text{index}/12-e_2/4-e_3/3-\text{cusps}/2$ gives zero in
each case. Thus even the asserted level-dividing-sixteen conclusion
would suffice. This proves reducibility.

For the characters, good supersingular reduction at $\ell$ is impossible:
its inertia action over the unramified base $\mathbb Q_\ell$ is
irreducible, whereas a global invariant line remains invariant locally.
In the good ordinary case, the inertial semisimplification is
$1\oplus\overline\chi_\ell$; the Tate and cusp cases have the same
semisimplification after their splitting twist. These are the local
ordinary/supersingular statements of \citet[Propositions 11--12]{AMSerre1972}.
Patch the inertia quadratic characters once more. Write the two global
characters as $\omega$ and $\overline\chi_\ell\omega^{-1}$. Both are
unramified away from $\ell$, and one is unramified at $\ell$. It is
therefore unramified at every finite prime and is trivial by Minkowski's
discriminant theorem. The two characters are the asserted ones.
Finally, the three Legendre two-torsion points are rational and a
quadratic twist acts on them by $[-1]=1$.
\end{proof}

\begin{statement}{A.7}{The modular cover of level $4\ell$}{Proved}
Let $Z_\ell$ be the finite \'etale cover of the tripod parametrizing a cyclic
order-$\ell$ subgroup of the Legendre curve, with $\ell$ odd. There is a
nonconstant $\mathbb Q$-morphism $Z_\ell\to Y_0(4\ell)$, finite after
smooth compactification. A connected component supporting an arithmetic
section is geometrically connected. The target has genus
$g(X_0(4\ell))=(\ell-3)/2$.
\end{statement}
\begin{proof}
Let $P=(0,0)$, $Q=(1,0)$, let $\phi:E\to E'=E/\langle P\rangle$
be the degree-two isogeny and $\phi':E'\to E$ its dual. The subgroup
$H_4=(\phi')^{-1}(\langle Q\rangle)$ is cyclic of order four: a lift
$R$ of $Q$ satisfies $2R=\phi(Q)\ne0$. If $C$ is the tautological
order-$\ell$ subgroup, $H_4+\phi(C)$ is cyclic of order $4\ell$ because
the orders are coprime. These base-change-compatible constructions give
the morphism. It is nonconstant: an elliptic curve has only finitely many
degree-two isogeny partners, whereas the Legendre $j$-invariant varies.
A nonconstant map between proper smooth curves is finite. The arithmetic
image of the fundamental group of a component is the Galois group of its
constant field; a section forces that field to be $\mathbb Q$.

For $N=4\ell$ the modular index is $6(\ell+1)$, the cusp number is
six, and both elliptic-point counts vanish. The genus formula gives
$1+6(\ell+1)/12-6/2=(\ell-3)/2$. The exact script checks the index,
cusp and elliptic-point counts at eight primes; the displayed calculation
proves the general formula. In particular level twenty has genus one.
\end{proof}

\begin{statement}{A.8}{The Selmer Fermat theorem at every prime at least five}{Proved}
For every prime $\ell\geq5$, every Selmer section of $U_\ell/\mathbb Q$
is cuspidal.
\end{statement}
\begin{proof}
The invariant line of A.6 gives a lift of the pushed-forward tripod
section to $Z_\ell$. It is Selmer by A.1. The target in A.7 has positive
genus, so \citet[Corollary 8.8]{AMStoll2007} gives its finite-abelian-descent
equality; finite-descent equality follows by inclusion. A.2 makes the
lifted section, hence its tripod image, a point or cusp section. In the
point case, the original section's containment in $\pi_1(U_\ell)$ fixes
a point of the finite \'etale fibre above the rational tripod point, so it
is induced by a rational point of $U_\ell$.

Its Legendre elliptic curve has full rational two-torsion. After the twist
of A.6 its stable order-$\ell$ line has character $1$ or
$\overline\chi_\ell$. In the former case it has a rational point of
order $\ell$. In the latter, quotient by its $\mu_\ell$ subgroup: the
dual isogeny's kernel is constant of order $\ell$. This quotient still
has full rational two-torsion, since an odd-degree isogeny is an
isomorphism on two-torsion. Thus in either case an elliptic curve over
$\mathbb Q$ has rational torsion containing
$\mathbb Z/2\oplus\mathbb Z/(2\ell)$, contrary to
\citet[Theorem (8)]{AMMazur1977}. This rules out the point case.

If the image is cuspidal, its lift lies in the stabilizer of a cusp in the
finite normalization fibre; that stabilizer is a cusp decomposition group
upstairs. Since the section surjects to $G_{\mathbb Q}$, the cusp is
rational. This completes the proof, including $\ell=5$.
\end{proof}

\begin{statement}{A.9}{The seven-torsion counterexample}{Refuted}
The assertion $Y_1(7)(\mathbb Q)=\varnothing$ in the last sentence of
\citet[Theorem 5.14]{AMSteklov2026} is false. The curve
$y^2-xy-4y=x^3-4x^2$ has discriminant $-1664$ and a point $(0,0)$ of
exact order seven. The correct exclusion needed in A.8 is simultaneous
full rational two-torsion and rational $\ell$-torsion for $\ell\geq5$.
\end{statement}
\begin{proof}
The successive nonzero multiples are
\[
 (0,0),\ (4,8),\ (2,2),\ (2,4),\ (4,0),\ (0,4),
\]
and the seventh multiple is the origin. The exact generalized
Weierstrass addition program checks each multiplication, each curve
equation and the nonzero discriminant. The final exclusion is exactly
Mazur's torsion list used in A.8, and does not assert emptiness of
$Y_1(7)$ alone.
\end{proof}

\begin{statement}{A.10}{The cubic}{Proved}
Every Selmer section of $U_3/\mathbb Q$ is cuspidal. Its three rational
cusps are all the rational points of the compactification.
\end{statement}
\begin{proof}
The rational functions
\[
 u=-\frac{3Z}{X+Y},\qquad
 v=\frac{9(X-Y)/(X+Y)-1}{2}
\]
give the elliptic equation $v^2+v=u^3-7$, since
\[
 (v^2+v-u^3+7)(X+Y)^3=27(X^3+Y^3+Z^3).
\]
The identity is verified symbolically in the exact program. The inverse
rational map exists from the ratios $Z/(X+Y)$ and $(X-Y)/(X+Y)$, so
smooth projective completion gives an isomorphism. The substitution
$(x,y)=(4u,8v+4)$ gives $y^2=x^3-432$, the model of $X_0(27)$ identified
in \citet[\S10.1]{AMKrizLi2019}.

Stoll's Corollary 8.8 gives finite-descent equality for this genus-one
modular curve. It also forces finiteness of its rational group: the
finite-descent locus is compact and closed, while an infinite finitely
generated group cannot be a compact countable topological group. For
completeness, a countable compact Hausdorff group has an isolated point by
the Baire theorem applied to its singleton decomposition; translations
make it discrete, hence finite. The diagonal rational group injects into
the finite adelic factors, so this argument applies to the asserted
equality.

At good primes $5$ and $7$ exact counting gives respectively $6$ and $9$
points. Prime-to-residue-characteristic torsion injects under reduction;
applying the two reductions prime by prime bounds rational torsion by
three. The points $O,(3,-5),(3,4)$ give all three and are the images of
the three rational coordinate cusps. A.2 therefore puts every Selmer
section into a rational cusp packet. No numerical $L$-value or unverified
analytic-rank table enters this argument.
\end{proof}

\begin{statement}{A.11}{Every odd exponent}{Proved}
For every odd $n\geq3$, every Selmer section of $U_n/\mathbb Q$ is
cuspidal. Its possible rational cusp centres are exactly
$[0:1:-1]$, $[1:0:-1]$, and $[1:-1:0]$.
\end{statement}
\begin{proof}
Choose a prime $\ell\mid n$. The morphism
\[
 [X:Y:Z]\longmapsto
 [X^{n/\ell}:Y^{n/\ell}:Z^{n/\ell}]
\]
is a finite \'etale cover $U_n\to U_\ell$ over $\mathbb Q$.
Functoriality carries local points to points and cusp sections to cusp
sections, so preserves Selmerness. By A.8 or A.10 the image is cuspidal.
The fixed-coset argument of A.8 in the finite normalization fibre puts
the original section into a rational cusp decomposition group. On a
coordinate line the ratio of the nonzero coordinates is a rational
$n$th root of $-1$; for odd $n$ it is $-1$. This gives exactly the three
displayed centres.

This is a theorem for Selmer sections. It has used no local section
conjecture and no assertion from the global section of the October 2026
preprint. The separate passage to arbitrary sections is treated below.
\end{proof}

\subsection{The quartic and every classical exponent}

\begin{statement}{A.12}{The three quartic quotients and their CM orders}{Proved}
Let $\overline C_4/\mathbb Q$ be $x^4+y^4=z^4$, and let $C_4$ denote
its complement of $xyz=0$. Its three degree-two quotient maps are
\[
 \begin{array}{c|c|c}
 (u,v)&\text{quotient equation}&\text{elliptic model}\\\hline
 (x/y,z^2/y^2)&v^2=1+u^4&E_-:Y^2=X^3-4X\\
 (x/z,y^2/z^2)&v^2=1-u^4&E_+:Y^2=X^3+4X\\
 (y/z,x^2/z^2)&v^2=1-u^4&E_+:Y^2=X^3+4X
 \end{array}
\]
All maps extend over the smooth projective completions. All three
geometric endomorphism rings are $\mathbb Z[i]$, with the nontrivial CM
endomorphisms defined over $\mathbb Q(i)$; their endomorphism rings over
$\mathbb Q$ are $\mathbb Z$. In particular none has complex
multiplication by $\mathbb Z[\sqrt{-2}]$. There is a $\mathbb Q$-isogeny
\[
 \operatorname{Jac}(\overline C_4)\longrightarrow E_-\times E_+^2.
\]
\end{statement}
\begin{proof}
The equations follow by dividing the Fermat equation by $y^4$ or by
$z^4$. More generally, for $v^2=1+a u^4$ the formulas
\[
 X=\frac{2(v+1)}{u^2},\qquad
 Y=\frac{4(v+1)}{u^3},\qquad
 u=\frac{2X}{Y},\quad v=\frac{2X^3}{Y^2}-1
\]
give inverse birational maps to $Y^2=X^3-4aX$. The exact identity is
checked in \texttt{a\_main\_quartic.py}. Birational maps between smooth
proper curves extend uniquely to isomorphisms. We choose $(u,v)=(0,1)$
as origin, which maps to the elliptic point at infinity. Each initial
map has degree two: the omitted coordinate is recovered by a quadratic
equation, and its sign is the nontrivial deck transformation.

On either elliptic model the map $(X,Y)\mapsto(-X,iY)$ fixes the origin
and squares to $[-1]$. Thus the geometric endomorphism ring contains
$\mathbb Z[i]$. In characteristic zero an elliptic endomorphism algebra
containing this field is that field, and $\mathbb Z[i]$ is its maximal
order. This proves equality. Complex conjugation sends $i$ to $-i$,
so its fixed endomorphism ring is $\mathbb Z$.

The three sign involutions on $\overline C_4$ form a Klein four group.
Its quotient is the conic $x^2+y^2=z^2$ in the squared coordinates,
which has genus zero. Hence there are no invariant differentials under
the whole group. Each quotient of degree two has genus one, so its
pullback differentials form a one-dimensional invariant subspace for
one of the three involutions. These are the three distinct nontrivial
characters of the Klein four group. In characteristic zero their sum
is direct, and it has dimension three, the genus of the smooth plane
quartic. The product of the induced Albanese maps has invertible
differential, hence finite kernel and equal-dimensional image. It is
the asserted isogeny, defined over $\mathbb Q$.
\end{proof}

\begin{statement}{A.13}{Exact Mordell--Weil groups of the quartic quotients}{Proved}
The rational groups of the elliptic curves in A.12 are
\[
 \begin{aligned}
 E_-(\mathbb Q)&=\{O,(0,0),(2,0),(-2,0)\}
                      \simeq(\mathbb Z/2)^2,\\
 E_+(\mathbb Q)&=\{O,(0,0),(2,4),(2,-4)\}
                      \simeq\mathbb Z/4.
 \end{aligned}
\]
Consequently $v^2=1+u^4$ has precisely $(0,1),(0,-1)$ and its two
rational points at infinity, whereas $v^2=1-u^4$ has precisely
$(0,1),(0,-1),(1,0),(-1,0)$.
\end{statement}
\begin{proof}
We give an exact two-isogeny descent, independent of a rank table.
Write $E_B:Y^2=X^3+BX$ and $E'_B=E_{-4B}$. The isogeny
\[
 \phi_B:E_B\longrightarrow E'_B,\qquad
 (X,Y)\longmapsto
 \left(X+\frac B X,\;Y\left(1-\frac B{X^2}\right)\right)
\]
has kernel $\{O,(0,0)\}$. It extends at these two points, and its dual
is the corresponding isogeny on $E'_B$ followed by the scaling
$(X,Y)\mapsto(X/4,Y/8)$. Their composition is multiplication by two.
The Kummer boundary for the dual isogeny is
\[
 \alpha_B:E_B(\mathbb Q)\longrightarrow
 \mathbb Q^\times/\mathbb Q^{\times2},\qquad
 \alpha_B(O)=1,\quad \alpha_B(0,0)=B,\quad
 \alpha_B(X,Y)=X\quad(X\ne0).
\]
Its kernel is the image of the dual isogeny. The formula is also checked
directly: an image has square $X$-coordinate, and if $X=r^2$ the
quadratic equation for a preimage has roots
$2r^2\pm2Y/r$, both rational. The exceptional points give the stated
values. The two isogeny exact sequences therefore yield
\[
 2^{\operatorname{rank}E_B(\mathbb Q)}
   =\frac{|\alpha_B(E_B(\mathbb Q))|
           |\alpha_{-4B}(E_{-4B}(\mathbb Q))|}{4}.
 \tag{A13.1}
\]
To see the factor four explicitly, compare $E_B/2E_B$ first with
$E_B/\phi_B' E'_B$, then with $E'_B/\phi_BE_B$; the intermediate
kernel is the rational kernel of $\phi_B'$, modulo its intersection
with $\phi_BE_B$. That intersection has order two exactly when
$E_B$ has full rational two-torsion. In the two cases its order
combines with $|E_B[2](\mathbb Q)|$ to give the same factor four.

Only square classes supported at primes dividing $B$ can occur:
at any other prime the equation forces $v_p(X)$ to be even.
For $B\in\{-4,4,16,-16\}$ the candidates are thus
$d\in\{1,-1,2,-2\}$. A nonexceptional point in class $d$ is equivalent
to a solution
\[
 N^2=dM^4+(B/d)e^4,\qquad \gcd(M,e)=1,
 \tag{A13.2}
\]
obtained by putting $X=dM^2/e^2$ and clearing denominators.
For $B>0$ every nonzero real $X$-coordinate is positive, excluding
the negative classes. For $B=16,d=2$, the right-hand residues of
(A13.2) modulo sixteen, with $M,e$ not both even, are
$\{2,8,10\}$; none is a square. For $B=-16,d=2$ they are the same,
and for $B=-16,d=-2$ they are $\{6,8,14\}$, again disjoint from
$\{0,1,4,9\}$. These exhaustive residue checks are in the program.
The evident points on the four curves now give exactly
\[
 \begin{array}{c|c}
 B&\alpha_B(E_B(\mathbb Q))\\\hline
 -4&\{1,-1,2,-2\}\\
 16&\{1\}\\
 4&\{1,2\}\\
 -16&\{1,-1\}.
 \end{array}
\]
Equation (A13.1) proves that both required ranks are zero.

Exact counting at the good primes three and five gives
$(|E_-(\mathbb F_3)|,|E_-(\mathbb F_5)|)=(4,4)$ and
$(|E_+(\mathbb F_3)|,|E_+(\mathbb F_5)|)=(4,8)$.
Prime-to-residue-characteristic torsion injects under reduction;
these two counts, applied prime by prime, bound either rational torsion
group by four. The displayed four points therefore exhaust each group.
The tangent formula gives $2(2,4)=(0,0)$ on $E_+$, proving its cyclic
structure. Finally the birational formulas of A.12 identify the stated
quartic points, including the two branches at infinity in the plus
quartic. This proves all assertions.
\end{proof}

\begin{statement}{A.14}{Exact nonvanishing and finite Tate--Shafarevich groups}{Proved}
Both elliptic curves $E_-$ and $E_+$ in A.12 have $L(E,1)>0$ and finite
Tate--Shafarevich group over $\mathbb Q$. The same assertions hold for
$E_{27}:v^2+v=u^3-7$. This nonvanishing is exact, not a numerical
analytic-rank assignment.
\end{statement}
\begin{proof}
The precise modular identifications in
\citet[Theorem 2]{AMMartinOno1997} are
\[
 \begin{aligned}
 f_{27}(\tau)&=\eta(3\tau)^2\eta(9\tau)^2,\\
 f_+(\tau)&=\eta(4\tau)^2\eta(8\tau)^2,\\
 f_-(\tau)&=\frac{\eta(8\tau)^8}
                    {\eta(4\tau)^2\eta(16\tau)^2},
 \end{aligned}
\]
with the indicated curves and levels $27,32,64$, respectively. Thus
the entire Hasse--Weil $L$-function, not a rounded special value, is
the Mellin transform of the corresponding form. For $m,t>0$ the
convergent product
\[
 \eta(imt)=e^{-\pi mt/12}
              \prod_{r\geq1}(1-e^{-2\pi mrt})
\]
is strictly positive. Every displayed quotient is consequently
positive on the positive imaginary axis. Since these are cusp forms,
their Mellin integrals converge at both endpoints; at zero this follows
by applying the modular transformation at that cusp. We obtain
\[
 L(E,1)=2\pi\int_0^\infty f_E(it)\,dt>0.
\]
This proves the needed nonvanishing without a denominator estimate
or a floating-point computation. As independent finite checks, our
program verifies the eta congruences, positive cusp orders, and the
coefficients against exact finite-field point counts through the
Sturm bounds eight and sixteen for the two Gaussian curves. These
checks support but do not replace the stated modular-identification
theorem.

The hypotheses of \citet[Theorem A, p.~527, and Remark (3), p.~528]
{AMRubin1987} now hold: the curves are defined over $\mathbb Q$, have
complex multiplication by an imaginary quadratic field, and have
nonzero $L$-value at one. Their Tate--Shafarevich groups over
$\mathbb Q$ are finite. The Gaussian CM was proved in A.12; for
$E_{27}$ the model $y^2=x^3-432$ has the automorphism
$(x,y)\mapsto(\zeta_3x,y)$ and therefore CM by the maximal order
$\mathbb Z[\zeta_3]$. This also gives an independent arithmetic
verification of the rank-zero finite-descent input in A.10, using
Rubin's Corollary C on p.~528 for the rank if desired.
\end{proof}

\begin{statement}{A.15}{Three separate proofs of the quartic Selmer theorem}{Proved}
Every Selmer section of the full profinite arithmetic fundamental group
of $C_4/\mathbb Q$ is cuspidal. Each of the three quotient maps in
A.12 proves this assertion separately. The rational cusp centres are
$[0:1:1]$, $[0:-1:1]$, $[1:0:1]$, and $[-1:0:1]$.
\end{statement}
\begin{proof}
For either elliptic target $E$, A.13 proves $E(\mathbb Q)$ finite and
A.14 proves $\operatorname{Sha}(E/\mathbb Q)$ finite. We spell out
the finite-descent implication, as in
\citet[Corollary 6.2 and its preceding argument]{AMStoll2007}.
The multiplication-by-$m$ isogenies form a cofinal family of geometric
finite \'etale covers of $E$. Their twisted descent sets have inverse
limit $\widehat{\operatorname{Sel}}(\mathbb Q,E)$, and Kummer theory gives
\[
 0\longrightarrow\widehat{E(\mathbb Q)}
  \longrightarrow\widehat{\operatorname{Sel}}(\mathbb Q,E)
  \longrightarrow T\operatorname{Sha}(E/\mathbb Q)
  \longrightarrow0.
\]
The last term vanishes because the group is finite. The first term is
$E(\mathbb Q)$ because that group is finite. Passing to local Kummer
maps, including connected components at the real place, gives
$E(\mathbb A_{\mathbb Q})_\bullet^{\mathrm{f\text{-}cov}}
=E(\mathbb Q)$.

Apply A.2 to any one of the three nonconstant proper maps
$\overline C_4\to E$. We obtain a point or cusp section over a rational
point of $\overline C_4$. The list of A.13 excludes an interior rational
point through each map separately. For the first quotient, its rational
points have $u=0$ or are at infinity, forcing $x=0$ or $y=0$. For the
second quotient, its rational points have $u=0$ or $v=0$, forcing $x=0$
or $y=0$; the third is symmetric. The coordinate equations give exactly
the four displayed points. Thus each argument ends in a rational cusp,
as claimed. This proof uses the Selmer hypothesis at every finite place;
the density-one clause of A.2 concerns finite support only.
\end{proof}

\begin{statement}{A.16}{Selmer sections for every classical exponent}{Proved}
Let $C_n/\mathbb Q$ be $x^n+y^n=z^n$ with $xyz=0$ removed, and let
$n\geq3$. The Selmer section classes of its full profinite fundamental
group are exactly the images of the rational cusp packets. Their
centres are
\[
 \begin{cases}
 [0:1:1],\ [1:0:1],\ [1:-1:0],&n\ \text{odd},\\
 [0:1:1],\ [0:-1:1],\ [1:0:1],\ [-1:0:1],&n\ \text{even}.
 \end{cases}
\]
This conclusion uses no $p$-adic local section conjecture. If every
section of $C_n$ has the local point-or-cusp property at every
completion, the same description holds for all section classes.
\end{statement}
\begin{proof}
For odd $n$, the isomorphism $[x:y:z]\mapsto[x:y:-z]$ identifies
$C_n$ with $U_n$, so A.11 applies. If $n$ is even and has an odd
prime factor $\ell$, the power map with exponent $n/\ell$ is a
finite \'etale morphism $C_n\to C_\ell$. If $n$ is a power of two,
then $4\mid n$ and the corresponding power map is a finite \'etale
morphism $C_n\to C_4$. These assertions follow on the coordinate
torus from characteristic-zero Kummer theory; their only ramification
on the smooth completions is at the removed coordinate divisors.
The image of a Selmer section is Selmer. A.11 or A.15 makes it
cuspidal, and the normalization-fibre argument of A.8 makes the
original section cuspidal over a rational cusp of $C_n$.

For odd $n$ the only rational root of $1$ or $-1$ of the indicated
order is $1$ or $-1$, giving the first list. For even $n$, the equations
at $x=0$ and $y=0$ give both signs, while $z=0$ has no real point.
This gives the second list. Conversely every section in a rational
cusp packet is locally cuspidal at every completion and hence Selmer.
The packet parametrization, its real-place signs, and the precise
additional hypothesis for the assertion about all sections are treated
in the subsequent local analysis; we make no density-one substitution
for that all-place hypothesis here.
\end{proof}
